% -*- coding: utf-8 -*-
% XeLaTeX 编译
\documentclass[a3paper, landscape, 11pt, twocolumn]{article}

% 引入宏包
\usepackage{fontspec}
\usepackage[UTF8]{ctex}
\usepackage{geometry}
\usepackage{listings}
\usepackage{xcolor}
\usepackage{enumitem}
\usepackage{tabularx}
\usepackage{makecell}
\usepackage{fancyhdr}
\usepackage{amssymb}
\usepackage{lastpage}

% 设置页面边距 (A3横向双栏)
\geometry{left=2cm, right=2cm, top=2.5cm, bottom=2.5cm, columnsep=1.8cm}

% 颜色定义
\definecolor{codebg}{RGB}{245,245,245}
\definecolor{keywordcolor}{RGB}{0,0,150}
\definecolor{commentcolor}{RGB}{0,128,0}
\definecolor{stringcolor}{RGB}{163,21,21}

% Java代码块样式设置
\lstdefinestyle{JavaExam}{
    language=Java,
    basicstyle=\ttfamily\small,
    backgroundcolor=\color{codebg},
    keywordstyle=\color{keywordcolor}\bfseries,
    commentstyle=\color{commentcolor},
    stringstyle=\color{stringcolor},
    showstringspaces=false,
    breaklines=true,
    frame=single,
    rulecolor=\color{black!50},
    tabsize=4,
    xleftmargin=1em,
    xrightmargin=1em,
    aboveskip=1em,
    belowskip=1em
}

% 列表样式
\newlist{options}{enumerate}{1}
\setlist[options]{label=\Alph*., nosep, leftmargin=*, itemsep=0.5em}

% 选择题选项命令
\newcommand{\chooseoption}{\begin{options}\item}

% 填空线命令
\newcommand{\blank}{\underline{\hspace{3cm}}}
\newcommand{\tfblank}{\underline{\hspace{0.8cm}}}

% 页眉页脚设置
\pagestyle{fancy}
\fancyhf{}
\fancyhead[C]{\Large\bfseries Java 程序设计模拟试卷（H卷）}
\fancyfoot[C]{第 \thepage\ 页 共 \pageref{LastPage}\ 页}
\fancyfoot[R]{得分：\underline{\hspace{2cm}}}
\fancyfoot[L]{考试时间：120分钟}
\renewcommand{\headrulewidth}{0.4pt}

\begin{document}

% 试卷装订线与密封线模拟
\twocolumn[{
	\begin{center}
		\vspace*{-1cm}
		{\huge\bfseries Java 程序设计模拟试卷（H卷）}\\[1em]
		{\large 闭卷 \hspace{2em} 满分：150分 \hspace{2em} 考试时间：120分钟}\\[1em]
		\renewcommand{\arraystretch}{1.5}
		\begin{tabular}{|c|c|c|c|c|c|c|}
			\hline
			题号  & 一 (40分)        & 二 (20分)        & 三 (20分)        & 四 (30分)        & 五 (40分)        & 总分 (150分)      \\
			\hline
			得分  & \hspace{1.5cm} & \hspace{1.5cm} & \hspace{1.5cm} & \hspace{1.5cm} & \hspace{1.5cm} & \hspace{1.5cm} \\
			\hline
			阅卷人 &                &                &                &                &                &                \\
			\hline
		\end{tabular}\\[1.5em]
		{\large 姓名：\underline{\hspace{3cm}} \hspace{2em} 学号：\underline{\hspace{4cm}} \hspace{2em} 班级：\underline{\hspace{3cm}}}
		\vspace{1em}
		\hrule
		\vspace{1.5em}
	\end{center}
}]

\section*{一、单项选择题（本大题共 20 小题，每小题 2 分，共 40 分）}
\textit{在每小题列出的四个备选项中只有一个是符合题目要求的，请将其代码填写在题后的括号内。错选、多选或未选均无分。}

\begin{enumerate}[label=\arabic*., itemsep=1em]
	\item Java 语言是由哪家公司开发的？
	      \begin{options}
		      \item Microsoft
		      \item Apple
		      \item Sun Microsystems
		      \item Google
	      \end{options}

	\item 下列哪个是 Java 中合法的 boolean 值？
	      \begin{options}
		      \item 0
		      \item false
		      \item "true"
		      \item TRUE
	      \end{options}

	\item 关于 `switch` 语句，以下说法正确的是？
	      \begin{options}
		      \item switch 语句中的表达式可以是 long 类型
		      \item case 后面必须跟 break 语句，否则会编译错误
		      \item switch 语句可以用于 String 类型（Java 7 及以上）
		      \item default 块必须放在所有 case 之后
	      \end{options}

	\item 下列哪个循环结构至少会执行一次循环体？
	      \begin{options}
		      \item for 循环
		      \item while 循环
		      \item do-while 循环
		      \item 增强 for 循环
	      \end{options}

	\item 关于数组的定义，正确的是？
	      \begin{options}
		      \item \texttt{int[] arr = new int[5]\{1,2,3,4,5\};}
		      \item \texttt{int arr[] = new int[5]\{1,2,3,4,5\};}
		      \item \texttt{int[] arr = \{1,2,3,4,5\};}
		      \item \texttt{int arr[] = \{1,2,3,4,5\};}
	      \end{options}

	\item 下列哪个选项可以正确定义一个有 3 行 4 列的二维数组？
	      \begin{options}
		      \item \texttt{int[][] arr = new int[3][4];}
		      \item \texttt{int arr[3][4] = new int[][];}
		      \item \texttt{int arr[][] = new int[][4];}
		      \item \texttt{int[] arr[] = new int[][];}
	      \end{options}

	\item 在 Java 中，所有类的直接或间接父类是？
	      \begin{options}
		      \item Object
		      \item Class
		      \item System
		      \item Root
	      \end{options}

	\item 下列关于 `this` 关键字的说法，正确的是？
	      \begin{options}
		      \item this 可以在静态方法中使用
		      \item this 可以调用本类的其他构造方法
		      \item this 指向当前类的父类对象
		      \item this 不能作为方法的返回值
	      \end{options}

	\item 关于方法重写（Override）的规则，错误的是？
	      \begin{options}
		      \item 重写方法的名称、参数列表必须与父类方法相同
		      \item 重写方法的返回值类型可以是父类方法返回值类型的子类
		      \item 重写方法的访问权限不能比父类方法更严格
		      \item 父类的私有方法可以被子类重写
	      \end{options}

	\item 下列哪个修饰符可以修饰局部变量？
	      \begin{options}
		      \item public
		      \item private
		      \item protected
		      \item final
	      \end{options}

	\item 关于接口，以下说法错误的是？
	      \begin{options}
		      \item 接口可以包含抽象方法
		      \item 接口可以包含默认方法和静态方法
		      \item 接口可以包含私有方法（Java 9 开始）
		      \item 接口可以包含实例变量
	      \end{options}

	\item 下列哪个类是所有异常类的父类？
	      \begin{options}
		      \item Error
		      \item Exception
		      \item Throwable
		      \item RuntimeException
	      \end{options}

	\item 在异常处理中，`finally` 块会在什么时候执行？
	      \begin{options}
		      \item 只有在 try 块没有异常时执行
		      \item 只有在 catch 块执行后执行
		      \item 无论是否发生异常都会执行（除非 JVM 退出）
		      \item 只有 try 块和 catch 块都执行后才执行
	      \end{options}

	\item 下列哪个集合类是基于哈希表实现的，并且不允许存储重复元素？
	      \begin{options}
		      \item ArrayList
		      \item LinkedList
		      \item HashSet
		      \item HashMap
	      \end{options}

	\item 关于泛型，以下说法正确的是？
	      \begin{options}
		      \item 泛型类型参数可以是基本数据类型
		      \item 泛型在运行时可以获取具体类型信息
		      \item 泛型可以用于类、接口、方法
		      \item 泛型不支持通配符
	      \end{options}

	\item 在 Java 中，实现 `Runnable` 接口与继承 `Thread` 类相比，优势是？
	      \begin{options}
		      \item 可以避免单继承的限制
		      \item 可以创建多个线程共享同一个目标对象
		      \item 代码更简洁
		      \item 以上都是
	      \end{options}

	\item 关于 `synchronized` 关键字，以下说法正确的是？
	      \begin{options}
		      \item 可以修饰类
		      \item 可以修饰构造方法
		      \item 可以修饰静态方法，锁定的是当前类的 Class 对象
		      \item 可以修饰成员变量
	      \end{options}

	\item 在 Java I/O 中，`File` 类的主要作用是？
	      \begin{options}
		      \item 读写文件内容
		      \item 管理文件和目录的元信息（如创建、删除、重命名）
		      \item 提供缓冲功能
		      \item 实现对象序列化
	      \end{options}

	\item 下列哪个类用于在 JDBC 中建立数据库连接？
	      \begin{options}
		      \item Statement
		      \item Connection
		      \item DriverManager
		      \item ResultSet
	      \end{options}

	\item 关于 Swing 中 `JButton` 的事件处理，以下说法正确的是？
	      \begin{options}
		      \item 单击按钮会触发 MouseEvent
		      \item 需要实现 ActionListener 接口并注册到按钮
		      \item 只能添加一个监听器
		      \item 按钮的默认布局管理器是 BorderLayout
	      \end{options}

	\item 在 Java 中，下列哪个关键字用于引入包中的类？
	      \begin{options}
		      \item package
		      \item import
		      \item include
		      \item using
	      \end{options}
\end{enumerate}

\section*{二、判断题（本大题共 10 小题，每小题 2 分，共 20 分）}
\textit{请判断下列说法的正误。正确的打"√"，错误的打"×"。}

\begin{enumerate}[label=\arabic*., itemsep=0.8em]
	\item Java 程序源文件名必须与 public 类名相同，且扩展名为 .java。（ \quad ）
	\item 在同一个类中，可以定义两个同名的方法，只要它们的参数列表不同即可。（ \quad ）
	\item 子类不能继承父类的私有成员。（ \quad ）
	\item 接口中的成员变量默认是 `public static final` 的。（ \quad ）
	\item `throw` 和 `throws` 关键字的作用相同，都是用来抛出异常。（ \quad ）
	\item `ArrayList` 和 `LinkedList` 都是线程安全的集合类。（ \quad ）
	\item 泛型方法可以在非泛型类中定义。（ \quad ）
	\item 调用 `Thread.sleep()` 方法会使当前线程暂停执行，但不会释放对象锁。（ \quad ）
	\item `FileReader` 和 `FileWriter` 是字节流类。（ \quad ）
	\item 在 UDP 通信中，`DatagramSocket` 负责发送和接收数据报，`DatagramPacket` 负责封装数据。（ \quad ）
\end{enumerate}

\section*{三、填空题（本大题共 5 小题，每空 2 分，共 20 分）}

\begin{enumerate}[label=\arabic*., itemsep=1em]
	\item Java 的数据类型分为两大类：\blank 类型和 \blank 类型。

	\item 面向对象程序设计的三个特征是封装、\blank 和 \blank。

	\item 在 Java 中，使用关键字 \blank 定义常量，使用关键字 \blank 定义抽象方法。

	\item 线程的优先级可以通过 \blank 方法设置，线程休眠可以使用 \blank 方法。

	\item 集合框架中，`Map` 接口的常用实现类包括 \blank 和 \blank。
\end{enumerate}

\section*{四、程序运行结果题（本大题共 5 小题，每小题 6 分，共 30 分）}
\textit{（阅读下列程序，写出运行后的输出结果）}

\begin{enumerate}[label=\arabic*., itemsep=1em]
	\item \textbf{写出下列程序的输出结果：}
	      \begin{lstlisting}[style=JavaExam]
public class Test1 {
    public static void main(String[] args) {
        int sum = 0;
        for (int i = 1; i <= 5; i++) {
            switch (i) {
                case 1:
                case 2: sum += i; break;
                case 3: sum += i * 2; break;
                default: sum += i;
            }
        }
        System.out.println(sum);
    }
}
\end{lstlisting}
	      \textbf{输出结果：}\underline{\hspace{8cm}}

	\item \textbf{写出下列程序的输出结果：}
	      \begin{lstlisting}[style=JavaExam]
class A { A() { System.out.print("A "); } }
class B extends A { B() { System.out.print("B "); } }
class C extends B { C() { System.out.print("C "); } }
public class Test2 {
    public static void main(String[] args) {
        C c = new C();
    }
}
\end{lstlisting}
	      \textbf{输出结果：}\underline{\hspace{8cm}}

	\item \textbf{写出下列程序的输出结果：}
	      \begin{lstlisting}[style=JavaExam]
public class Test3 {
    public static void main(String[] args) {
        Integer a = 127;
        Integer b = 127;
        Integer c = 128;
        Integer d = 128;
        System.out.println(a == b);
        System.out.println(c == d);
    }
}
\end{lstlisting}
	      \textbf{输出结果：}\underline{\hspace{8cm}}

	\item \textbf{写出下列程序的输出结果：}
	      \begin{lstlisting}[style=JavaExam]
public class Test4 {
    public static void main(String[] args) {
        int[] arr = {1, 2, 3, 4, 5};
        try {
            System.out.println(arr[5]);
        } catch (ArrayIndexOutOfBoundsException e) {
            System.out.println("索引越界");
        } catch (Exception e) {
            System.out.println("其他异常");
        } finally {
            System.out.println("finally块");
        }
        System.out.println("程序结束");
    }
}
\end{lstlisting}
	      \textbf{输出结果：}\underline{\hspace{8cm}}

	\item \textbf{写出下列程序的输出结果：}
	      \begin{lstlisting}[style=JavaExam]
class Counter {
    static int count = 0;
    public Counter() { count++; }
    public static int getCount() { return count; }
}
public class Test5 {
    public static void main(String[] args) {
        Counter c1 = new Counter();
        Counter c2 = new Counter();
        Counter c3 = new Counter();
        System.out.println(Counter.getCount());
    }
}
\end{lstlisting}
	      \textbf{输出结果：}\underline{\hspace{8cm}}
\end{enumerate}

\newpage
\section*{五、编程题（本大题共 2 小题，每小题 20 分，共 40 分）}
\textit{（要求代码完整、结构清晰、注释合理）}

\begin{enumerate}[label=\arabic*.]

	\item \textbf{设计一个"动物"类层次结构（20分）\\
	      定义抽象类 `Animal`，包含属性 `name`（名称）和 `age`（年龄），提供构造方法初始化这些属性。定义抽象方法 `void makeSound()`。\\
	      派生两个子类：`Dog` 和 `Cat`。\\
	      - `Dog` 类增加属性 `breed`（品种），重写 `makeSound()` 方法输出"汪汪"。\\
	      - `Cat` 类增加属性 `color`（颜色），重写 `makeSound()` 方法输出"喵喵"。\\
	      在测试类 `Main` 中，创建一个 `ArrayList<Animal>`，添加若干 Dog 和 Cat 对象。遍历列表，对每个动物调用 `makeSound()` 方法，并输出动物的名称、年龄和特有属性（需进行类型判断）。}
	\item \textbf{实现一个简单的"图书馆"管理系统（20分）\\
	      定义 `Book` 类，属性包括 `id`（编号，String）、`title`（书名）、`author`（作者）、`isBorrowed`（是否借出，boolean）。提供构造方法、getter/setter，以及 `toString()` 方法。\\
	      定义 `Library` 类，使用 `HashMap<String, Book>` 存储图书（键为编号）。提供以下方法：\\
	      - `void addBook(Book book)`：添加图书，如果编号已存在则输出提示。\\
	      - `void removeBook(String id)`：根据编号删除图书。\\
	      - `Book searchBook(String id)`：根据编号查找图书。\\
	      - `void borrowBook(String id)`：借出图书（设置 isBorrowed = true），如果图书不存在或已借出则输出相应提示。\\
	      - `void returnBook(String id)`：归还图书（设置 isBorrowed = false），如果图书不存在或未借出则输出相应提示。\\
	      - `void displayAllBooks()`：显示所有图书信息。\\
	      在测试类中，添加至少 5 本图书，测试借书、还书、查找、删除等功能。}
\end{enumerate}

\newpage
\begin{center}
	{\large\bfseries —— 试卷结束，祝考试顺利 ——}
\end{center}

\end{document}
